\section{Evidence Availability and Reporting}\label{sec:reporting}
Attack adaptation and original-task utility occur in different mapped protection settings. Table~\ref{tab:availability} identifies four eligible query/workflow privacy comparisons: two document adaptation and one documents matched original-task utility; none locates both. The adapted settings share one preprint, while the utility observation is WRA's published visibility comparison.

\begin{table*}[t]\centering\small
\caption{Evidence located in the retained records. Counts use eligible settings; the second column also gives distinct sources. ``Not located'' leaves evidence unresolved within the inspected scope. These counts describe availability, not study quality or checklist validity.}\label{tab:availability}
\setlength{\tabcolsep}{4pt}
\begin{tabularx}{\textwidth}{@{}P{.28\textwidth}P{.15\textwidth}P{.11\textwidth}Y@{}}\toprule
Question & Eligible settings / sources & Located & Supporting settings and consequence \\\midrule
Protected-condition attack adaptation & 4 / 2 & 2/4 & Communication graphs: transformed captures and live OHTTP. Both come from one preprint. \\\addlinespace[3pt]
Original-task utility under the same intervention & 4 / 2 & 1/4 & WRA visibility comparison. It does not price WRA's executed-decoy condition. \\\addlinespace[3pt]
Adaptation and original-task utility in one setting & 4 / 2 & 0/4 & Adaptation and utility are located in different records; transport cost alone does not fill the task-outcome field. \\\addlinespace[3pt]
Original-task utility for executed cover & 1 / 1 & 0/1 & WRA's five research decoys report recovery and extra research; matched original-task utility is not located. \\\bottomrule
\end{tabularx}
\end{table*}


\paragraph{Different objectives need different denominators}
Three identity-obfuscation settings, two adversarial-Agent blocking settings and thirty-five without a protection comparison remain outside this denominator. The access map separately counts twenty-two transit settings, ten recipient settings, eleven active settings and one downstream preference analysis. These author-assigned classes partition settings; source groups overlap. The counts describe the retained corpus.

\paragraph{A compact reporting contract}
Arp et al. and WF guidelines establish sampling, labels, observation, splits and adaptive evaluation~\cite{arp2022dos,juarez2014critical,mathews2023critical}. Our eight-field checklist adds service/access mapping, label stage and original-task costs. Supplement S5 provides the checkable form, guideline correspondence, eligibility ledger and completed WRA, OHTTP and ANT examples. Each entry records evidence, a missing item or inapplicability.

The reporting unit is the intervention: which observer remains, how the attack adapts, and what happens to the original task. Recording these together would close the mapped privacy--utility gap.
